\section{Preliminaries}\label{sec:prelim}

\subsection{The statement (T)}\label{sec:prelim-T}

We use the machine model of the first paper of this series~\cite{FirstPaper2026}: a deterministic Turing machine with
a read-only input tape, a constant number of work tapes over a finite alphabet, and a write-only output tape. For
$d\ge1$ let $b(d)=\max(1,\clog d)$. The input is a $d\times d$ matrix with $b(d)$-bit entries written row by row, of
length $L=d^2b(d)$, and the machine must write its transpose in the same encoding. Statement~(T) says that some
machine does this within $c(L+1)$ steps, for a constant $c$ and all $d$. Its formal version,
\lean{Transpose.LinearTimeTransposable}, is fixed in a frozen file of the first paper's library. A machine is
\emph{oblivious} if its head positions, output times and halting time depend only on the input length.

The Lean library of this paper does not model machines. Its theorems are about recovery relations, polymatroids,
storage codes and index codes on finite sets. The connection to~(T) goes through the fourth paper of the
series~\cite{FourthPaper2026} and through \emph{Directed Recovery Structures in Tape
Computations}~\cite{DirectedRecovery2026}; we recall it in \cref{sec:construction}.

\subsection{Recovery relations}\label{sec:prelim-relations}

A \emph{recovery relation} on a finite set $V$ is a map $N\colon V\to2^V$. We write $u\to v$, an \emph{arc}, when
$u\in N(v)$: then $u$ is a \emph{parent} (in-neighbour) of $v$ and $v$ a \emph{child} of $u$. The \emph{in-degree}
of $v$ is $\card{N(v)}$. The relation is \emph{loopless} if $v\notin N(v)$ for every $v$. For $K\subseteq V$, the
\emph{induced relation} $N[K]$ on $K$ is $v\mapsto N(v)\cap K$, and $N-Z=N[V\setminus Z]$. For a simple graph $G$ on
$V$, its \emph{neighbourhood relation} $N_G$ maps $v$ to its set of neighbours (\lean{MotionDepthCodes.nbhd}); every
arc of $N_G$ is reciprocated.

A \emph{digon} is a pair of arcs $u\to v$ and $v\to u$; an arc that is not part of a digon is \emph{one-way}. The
\emph{digon graph} $G_2(N)$ is the simple graph whose edges are the digons. The \emph{underlying graph} $U(N)$
(\lean{und}) joins $u\ne v$ whenever $u\to v$ or $v\to u$. A \emph{transitive triangle} is a triple of distinct
vertices $x,y,z$ with $x\to y$, $y\to z$ and $x\to z$; we call $y$ its \emph{middle vertex}. A relation is
\emph{TT-free} if it has no transitive triangle (\lean{TTFree}); \lean{TTFreeOn} is the same condition for triples
inside a set. In a loopless TT-free relation the parents of a vertex are pairwise non-adjacent, the digon graph is
triangle-free, and every triangle of the underlying graph is a \emph{directed $3$-cycle} $x\to y\to z\to x$ without
reverse arcs (\cref{lem:dtri} below); bidirected triangles are excluded.

A set $A\subseteq V$ is \emph{acyclic} for $N$ if every nonempty $T\subseteq A$ contains a vertex with no parent in
$T$ (\lean{IsAcyclicIn}); equivalently, $N[A]$ has no loops and no directed cycles. A relation is
\emph{$2$-dicolourable} if $V$ is the union of two disjoint acyclic sets.

\paragraph{Paired incidence.} \emph{Directed Recovery Structures}~\cite[Definition~2.3]{DirectedRecovery2026} attaches
to $N$ its incidence graph $E(N)$: a bipartite graph with vertices $L_p$ and $R_p$ for $p\in V$ and an edge
$L_pR_q$ whenever $p\in N(q)$. It puts $b_E(r)=\card{\{p\in V:\dist_E(L_p,R_p)\le r\}}$, and a sequence of relations
satisfies the \emph{paired-distance condition} if $b_E(r)=o(\card V)$ for every fixed $r$. Distances between $L_p$
and $R_p$ are odd. Our hypothesis ``TT-free'' is the exact form of this condition at radius three.

\begin{lemma}[Transitive triangles are paired paths of length three]\label{lem:tt-paired}
\notformalized
Let $N$ be a recovery relation and $p\in V$. Then $\dist_E(L_p,R_p)=1$ if and only if $p\in N(p)$. If $N$ is
loopless, then $\dist_E(L_p,R_p)=3$ if and only if $p$ is the middle vertex of a transitive triangle. In particular,
a loopless relation is TT-free if and only if $b_E(3)=0$, that is, if and only if $\dist_E(L_p,R_p)\ge5$ for every
$p$.
\end{lemma}

\begin{informalproof}
The edge $L_pR_p$ exists exactly when $p\in N(p)$. Let $N$ be loopless. A path $L_p,R_z,L_x,R_p$ of length three
means $p\in N(z)$, $x\in N(z)$ and $x\in N(p)$, that is, $p\to z$, $x\to z$ and $x\to p$. Its vertices are distinct,
so $z\ne p$ and $x\ne p$, and $x\ne z$ because $x\in N(z)$ and $N$ is loopless. Hence $x\to p\to z$ with $x\to z$ is
a transitive triangle with middle vertex $p$. Conversely, every transitive triangle $x\to p\to z$, $x\to z$ gives
this path. Since the distance is odd and not $1$, it is at least $5$ when no such path exists.
\end{informalproof}

\subsection{Polymatroids and storage codes}\label{sec:prelim-codes}

We use the definitions of the fourth paper~\cite[Definitions~2.1 and~2.2]{FourthPaper2026}. A \emph{polymatroid} on
$V$ is a function $h\colon2^V\to\R$ with $h(\emptyset)=0$ that is monotone and submodular. It is \emph{valid} for
$N$ if $h(N(v)\cup\{v\})=h(N(v))$ for every $v\in V$. We write $h(v)$ for $h(\{v\})$,
\[
  h(A\mid B)=h(A\cup B)-h(B),\qquad I(A;B\mid S)=h(A\cup S)+h(B\cup S)-h(A\cup B\cup S)-h(S)\ge0,
\]
and $I(w;X)=h(w)+h(X)-h(X\cup\{w\})$ (\lean{mi}). The subject of this paper is the \emph{weighted half bound}
\begin{equation}\label{eq:star}
  2h(V)\;\le\;\sum_{v\in V}h(v)\qquad\text{for every polymatroid $h$ valid for $N$.}\tag{$\star$}
\end{equation}
In the terms of the fourth paper, \eqref{eq:star} says that $N$ satisfies the rate bound $\tfrac12$; the
formal statement is \lean{MotionDepthCodes.RateBound}~$N\,(1/2)$. A \emph{storage code} for $N$ over a finite alphabet $\alpha$ is a set
$C\subseteq\alpha^V$ such that two words of $C$ that agree on $N(v)$ agree at $v$, for every $v$
(\lean{MotionDepthCodes.IsStorageCode}). The \emph{guessing number} $\gn(N,q)$ is the largest value of
$\log_q\card C$ over storage codes with $\card\alpha=q$, and $\gn(N)=\sup_q\gn(N,q)$ is the asymptotic guessing
number~\cite{Riis2007,ChristofidesMarkstrom2011}.

\begin{theorem}[Rate bounds bound codes]\label{thm:codes}
\cited{the fourth paper~\cite[Theorem~6.3 and Corollary~6.9]{FourthPaper2026}, Lean-verified there
(\lean{MotionDepthCodes.card_sq_le_of_rateBound_half}).}
If $N$ satisfies \eqref{eq:star}, then every storage code $C\subseteq\alpha^V$ for $N$ over a nonempty finite
alphabet satisfies $\card C^2\le\card\alpha^{\card V}$.
\end{theorem}

The proof applies \eqref{eq:star} to the entropy of a uniformly random codeword, which is a valid
polymatroid. The same argument gives $\card C^2\le\prod_v\card{\alpha_v}$ for codes with coordinate alphabets
$\alpha_v$~\cite[Corollary~3.4]{CircularMemo2026}; we do not use this.

\begin{lemma}[Closure]\label{lem:closure}
\cited{the fourth paper~\cite[Lemma~6.4]{FourthPaper2026}, Lean-verified there
(\lean{MotionDepthCodes.IsPolymatroid.eq_of_closure}).}
Let $h$ be a polymatroid valid for $N$, and let $S\subseteq T$ be such that $N(v)\subseteq S$ for every
$v\in T\setminus S$. Then $h(T)=h(S)$.
\end{lemma}

\begin{lemma}[Fewer arcs]\label{lem:fewer-arcs}
\cited{the fourth paper's Theorem~H~\cite[Theorem~6.5]{FourthPaper2026}, applied to the identity map; Lean-verified
there (\lean{MotionDepthCodes.RateBound.of_hom}).}
If $N(v)\subseteq N'(v)$ for every $v$ and $N'$ satisfies \eqref{eq:star}, then so does $N$.
\end{lemma}

Indeed, a polymatroid valid for $N$ is valid for $N'$, by diminishing returns. The last tool is contraction.

\begin{lemma}[Contraction]\label{lem:contraction}
\formal{Decomposition.lean}{contr, contr_isPolymatroid, contr_insert_eq}
Let $h$ be a polymatroid on $V$ and $U\subseteq V$. Then $h_U(X)=h(X\cup U)-h(U)$ is a polymatroid. If $h$ is valid
for $N$ and $N(v)\subseteq U\cup X$, then $h_U(X\cup\{v\})=h_U(X)$; in particular the restriction of $h_U$ to
$V\setminus U$ is valid for $N-U$. Moreover $h_U(w)=h(w)-I(w;U)$ for $w\notin U$.
\end{lemma}

\begin{outline}
The polymatroid axioms for $h_U$ follow from those for $h$, because $X\mapsto X\cup U$ preserves unions,
intersections and inclusions. If $N(v)\subseteq U\cup X$, the closure lemma (\cref{lem:closure}) gives
$h(X\cup U\cup\{v\})=h(X\cup U)$. The last identity is the definition of $I(w;U)$.
\end{outline}

\subsection{Formal status of the results}\label{sec:prelim-lean}

Every numbered result names its formal status in one of four ways. A line
\begin{center}
{\small\textsf{Lean:} \texttt{declaration, \dots}\hfill[\texttt{file}]}
\end{center}
gives the Lean declarations of a fully formalized result. The declarations live in the namespace \texttt{HalfRate}
and the files in the directory \leanfile{lean/HalfRate/}. The text after such a statement, labelled \emph{Outline of the
formal proof}, is a guide to the formal proof; it is a convenience, not a proof. A line beginning \textsf{Lean
(partial)} marks a result whose formal counterpart covers all but an explicitly named part, which the text proves.
The words \textsf{Not formalized} mark a result with no formal counterpart; it comes with a complete proof, labelled
\emph{Proof (not formalized)}. A line beginning \textsf{Cited} marks a result of an earlier paper of the series that is
used without proof, and says where it is proved and whether it is Lean-verified there. Where the English statement
simplifies, the Lean statement is authoritative. Numerical observations are labelled as such and are not used in
proofs. Numbered definitions carry a status line only when they are formalized.

The library depends on Lean~4 and Mathlib (both v4.34.1), on the Lean library of the fourth paper (whose
polymatroids, storage codes and rate bounds we use unchanged) and on that of its circular memo~\cite{CircularMemo2026},
imported through relative paths. It contains no \texttt{sorry}, no \texttt{native\_decide} and no unsafe evaluation,
and every declaration cited here depends only on the axioms \texttt{propext}, \texttt{Classical.choice} and
\texttt{Quot.sound} (\cref{app:verify}). \Cref{app:index} lists every declaration of this library cited in the paper.
